A distinctive consequence of basing \Gweek{} on CBPV is its treatment of
\texttt{let} bindings. In the core calculus, a \texttt{let} binding corresponds
to the sequencing construct $\CTerm{\CTo{M}{x}{N}}$---the monadic bind. In a
strict evaluator, $\CTerm{M}$ would be fully evaluated before control passes to
$\CTerm{N}$. But in a language with nondeterminism, this is not the only
reasonable choice, and the difference turns out to be consequential.

\subsection{The Mechanism}

The key transition is the bind rule in \cref{figure:machine}. In a standard CK
machine, $\CTerm{\CTo{M}{x}{N}}$ pushes a continuation onto the stack and
evaluates $\CTerm{M}$ immediately. \Gweek{} instead \emph{suspends}
$\CTerm{M}$: it records the computation in the suspension environment
$\Susp{S}$ and continues with $\CTerm{N}$:
%
\[
  \SLMachine{\CTo{M}{x}{N}}{K}
  \quad\MStep\quad
  \SLMachine{N}[\AddSusp{S}{x}{M}]{K}
\]
%
The computation $\CTerm{M}$ is not evaluated. The suspension environment is
extended with the binding $x \mapsto \CTerm{M}$, and the variable $x$ in
$\CTerm{N}$ refers to this suspended computation.

The suspension is \emph{forced}---the deferred computation is actually
evaluated---only when $x$ is encountered in a position that \emph{demands} its
value: a case expression whose scrutinee is $x$, or a unification whose operand
is $x$. When this happens, the machine reverses the suspension: it extracts the
suspended computation from $\Susp{S}$, pushes a special \emph{set} frame onto
the stack that records which suspension to update, and begins evaluating the
deferred computation:
%
\[
  \SLMachine{M}[\DisjSusp{S}{x}{N}]{K}
  \quad\MStep\quad
  \SLMachine{N}{\StkKont{x}{M}{K}}
  \qquad\text{(when $x$ is needed in $\CTerm{M}$)}
\]
%
When the deferred computation completes with a value, the set frame updates the
suspension environment so that subsequent references to $x$ resolve immediately.
The effect is exactly what the strict machine would have done---but deferred
until the point of use.

If $x$ is never needed during normal execution---if it is used only in
positions that do not inspect its value, such as being passed as an argument or
stored in a data structure---the suspension remains unforced until the machine
reaches a terminal state (a return with an empty stack). At that point the
machine enters a \emph{drain phase}: it iterates over the remaining pending
suspensions, forcing each in turn by pushing a set frame and evaluating the
deferred computation. This ensures that any effects the suspensions contain (in
particular, failures that should prune the current branch) are properly
accounted for before a result is reported.

\subsection{Why Suspensions Matter}

The difference between the strict and lazy transitions may seem minor, but it
has a profound effect on the interaction between sequencing and nondeterminism.

Consider a program that searches for a sorted permutation of a list:
%
\begin{haskellcode}
let ys = perm [3, 1, 2] in
let c = guard_sorted ys in
ys.
\end{haskellcode}
%
\noindent Under the strict semantics, the machine evaluates
\mintinline{haskell}|perm [3, 1, 2]| first. This is a nondeterministic
computation that branches into all $3! = 6$ permutations. Each branch is a
separate machine state. Only then does the machine evaluate
\mintinline{haskell}|guard_sorted ys| in each branch independently. The
constraint check cannot prune branches that were created \emph{before} it ran.

Under the suspension semantics, the evaluation proceeds differently. First,
\mintinline{haskell}|perm [3, 1, 2]| is suspended---no branches are created
yet. Then \mintinline{haskell}|guard_sorted ys| is also suspended. The machine
reaches the result expression \mintinline{haskell}|ys|, which is a suspension
reference. If the context demands the value of \mintinline{haskell}|ys|---for
example, to print the result---then \mintinline{haskell}|ys| is forced. The
permutation is generated incrementally: each element is produced and immediately
available to the sortedness check. When \mintinline{haskell}|guard_sorted|
examines the first two elements and finds them unsorted, it fails, pruning the
branch \emph{before} the remaining elements are ever chosen.

The equational theory illuminates this. The
\emph{bind-return law} states that $\CTerm{\CTo{\Return{V}}{x}{M}} =
\CTerm{\CSubst{M}{V}{x}}$: when the bound computation simply returns a value,
sequencing is equivalent to substitution. But when $\CTerm{M}$ is
nondeterministic---when it branches into multiple alternatives---the bind
genuinely sequences, and the continuation $\CTerm{N}$ runs in each branch. The
suspension mechanism defers this branching, allowing the continuation to proceed
and potentially constrain the search before the full tree is materialised.

\subsection{Strict Mode}

\Gweek{} supports a strict evaluation mode (the \texttt{--strict} flag) in
which every \texttt{let} binding uses the standard CK transition instead of the
suspension transition. The machine pushes a continuation frame and evaluates the
bound computation immediately.

Strict mode is appropriate for programs whose search space is inherently
finite, or where all constraints are forced immediately by case expressions. It
eliminates the overhead of the suspension environment and avoids the subtleties
of deferred evaluation.

However, strict mode changes the termination behaviour of the system. Programs
that rely on lazy interleaving of generation and testing may fail to terminate
in strict mode, because the generation fully unfolds before any constraint can
prune it. The spectrum between lazy and strict FLP semantics---and the question
of which equations hold in each---is an interesting direction for future work.

\subsection{Programming Discipline}

The suspension mechanism requires the programmer to attend to two patterns.

\paragraph{Forcing constraints}
If a constraint is bound to a variable that is never scrutinised, the
constraint is suspended and never forced during execution. For example, in
\mintinline{haskell}|let c = neq a b in body|, if \mintinline{haskell}|c| is
never case-analysed in \mintinline{haskell}|body|, the call to
\mintinline{haskell}|neq| has no effect. The fix is to force the result:
%
\begin{haskellcode}
case neq a b of
    Z -> body
  | S n -> fail.
\end{haskellcode}
%
\noindent The case scrutinee demands the head constructor, which forces the
suspension and triggers any embedded failure.

\paragraph{Interleaving generation and testing}
Even when constraints are properly forced, the order in which nondeterministic
choices and constraints are combined matters for efficiency. Generating all
candidates monolithically and then testing each means the full nondeterministic
tree is explored before any constraint can prune it. The efficient alternative
is to \emph{interleave} generation and testing: produce one candidate element,
check applicable constraints, and only then produce the next.
